% from ega3-1-fr.tex
\oldpage[III]{114}
Idéaux cohérents $\mathcal J\mathcal B$ et
$\mathcal J\mathcal K=\mathcal J(\mathcal K\mathcal B)
=\mathcal K(\mathcal J\mathcal B)$, de sorte que l'on a le diagramme
d'inclusions
\[
\xymatrix@R=1.5em{
\mathcal J\mathcal B\ar@{^{(}->}[r] & \mathcal B\\
\mathcal J\ar@{^{(}->}[r]\ar@{^{(}->}[u] &
\mathcal A\ar@{^{(}->}[u]\\
\mathcal J\mathcal K\mathcal B=\mathcal J\mathcal K
\ar@{^{(}->}[r]\ar@{^{(}->}[u] &
\mathcal K\ar@{^{(}->}[u]
}
\tag{2.6.2.6}\label{III.2.6.2.6-fr}
\]
Soit $\mathcal J'$ l'Idéal cohérent $(\mathcal J\mathcal B)^\sim$ de
$\mathcal O_{X'}$, de sorte que $\mathcal J\mathcal B=g_*(\mathcal J')$,
$\mathcal J'\mathcal K'=(\mathcal J\mathcal K\mathcal B)^\sim$, et par
suite
\[
\mathcal J'/\mathcal J'\mathcal K'
=(\mathcal J\mathcal B/\mathcal J\mathcal K\mathcal B)^\sim
\]
(II, 1.4.4). Comme $\mathcal J|V=\mathcal J\mathcal K|V$ pour tout
ouvert $V$ ne rencontrant pas $Z$, on voit que le support de
$\mathcal J'/\mathcal J'\mathcal K'$ est contenu dans $Z'$. Comme $Z'$ est
propre sur $Y$, on peut appliquer (\hyperref[III.2.2.4-fr]{2.2.4}) et on
voit que pour $n$ assez grand, l'application canonique
\[
\Gamma(X',\mathcal J'\otimes\mathcal L'^{\otimes n})
\longrightarrow
\Gamma(X',(\mathcal J'/\mathcal J'\mathcal K')
\otimes\mathcal L'^{\otimes n})
\]
est surjective. Mais en vertu de (0\textsubscript{I}, 5.4.10), on en
conclut que l'application canonique
\[
\Gamma(X,\mathcal J\mathcal B\otimes\mathcal L^{\otimes n})
\longrightarrow
\Gamma(X,(\mathcal J\mathcal B/\mathcal J\mathcal K\mathcal B)
\otimes\mathcal L^{\otimes n})
\]
est surjective.

Cela étant, soient $i:Z\to X$ l'injection canonique,
$i':Z'\to X'$ l'injection canonique, de sorte qu'on a le diagramme
commutatif
\[
\xymatrix{
X'\ar[d]_g & Z'\ar[l]_{i'}\ar[d]^{h}\\
X & Z\ar[l]^{i}
}
\]
Soient $\mathcal M=i^*(\mathcal L)$, $\mathcal M'=i'^*(\mathcal L')$~;
comme $\mathcal L'$ est ample, $\mathcal M'$ est ample (II, 5.1.12), et
d'autre part $\mathcal M'=h^*(\mathcal M)$~; on conclut donc de
l'hypothèse de récurrence noethérienne (puisque $Z\neq X$) que
$\mathcal M$ est ample. Par suite
$i^*(\mathcal J/\mathcal J\mathcal K)\otimes\mathcal M^{\otimes n}$ est
engendré par ses sections au-dessus de $Z$ pour $n$ assez grand
(II, 4.5.5). Comme
$\mathcal J/\mathcal J\mathcal K=i_*(i^*(\mathcal J/\mathcal J\mathcal K))$,
on déduit encore de (0\textsubscript{I}, 5.4.10) qu'il existe une section
$s$ de $(\mathcal J/\mathcal J\mathcal K)\otimes
\mathcal L^{\otimes n}$ au-dessus de $X$ (pour un $n$ assez grand) telle
que $s(x)\neq0$, puisque l'on a $\mathcal J_x=\mathcal A_x$ par
définition de $\mathcal J$ et $\mathcal K_x\neq\mathcal A_x$ par
hypothèse. Le diagramme (\hyperref[III.2.6.2.6-fr]{2.6.2.6}) montre que
$s$ est aussi une section de
$(\mathcal J\mathcal B/\mathcal J\mathcal K\mathcal B)
\otimes\mathcal L^{\otimes n}$ au-dessus de $X$, donc $s$ est l'image
canonique d'une section $t$ de
$(\mathcal J\mathcal B)\otimes\mathcal L^{\otimes n}$ au-dessus de $X$.
Mais par définition, l'image canonique $s$ de $t$ mod.
$(\mathcal J\mathcal K\mathcal B)\otimes\mathcal L^{\otimes n}$ est
dans
\[
(\mathcal J\otimes\mathcal L^{\otimes n})/
(\mathcal J\mathcal K\mathcal B\otimes\mathcal L^{\otimes n})
\]
donc, en vertu de (\hyperref[III.2.6.2.6-fr]{2.6.2.6}), cela implique que
$t$ est une section de $\mathcal J\otimes\mathcal L^{\otimes n}$
au-dessus de $X$, et \emph{a fortiori} une section de
$\mathcal L^{\otimes n}$. On a vu ci-dessus que $t(x)\neq0$, donc
$x\in X_t$, et par définition de $\mathcal J$, $t(y)=0$ dans $X-U$ qui
est le support de $\mathcal O_X/\mathcal J$~; donc $X_t\subset U$, ce qui
achève la démonstration.
\end{env}
\end{proof}

